\section{Реализация программной системы}
\label{sec:impl}

\subsection{Общие принципы реализации}

Компоненты Master и Agent реализованы как ASP.NET Core Minimal API на .NET~8. Библиотека Chronos.Core публикуется как отдельный пакет и переиспользуется в консольных сценариях, тестах и серверах. Конфигурация задаётся переменными окружения и секциями \texttt{appsettings}, что упрощает контейнерный деплой.

\subsection{Точка входа Master}

В \texttt{Program.cs} мастер-узла настраиваются: подключение к PostgreSQL и автоматические миграции EF Core, регистрация \texttt{DbContextFactory}, HTTP-клиенты для вызовов агентов, endpoint-ы Swagger, статические файлы UI, фоновые службы (очистка устаревших агентов и аудита, репликация томов при наличии S3, генерация HAProxy при включённом каталоге динамической конфигурации).

\subsection{Точка входа Agent}

Агент регистрирует маршруты из \texttt{AgentRoutes} и \texttt{AgentMainRoutes}, поднимает планировщик проверок, при необходимости --- встроенный MinIO или внешнее объектное хранилище для сценария прямой выгрузки снимков. Параметры таймаутов и параллелизма тестов задаются конфигурацией, чтобы не перегружать слабые стенды.

\subsection{Реализация подсистемы TCP-маршрутизации}

Реестр маршрутов сохраняет JSON и генерирует фрагмент \texttt{chronos-tcp.cfg}. Запись выполняется во временный файл с последующим переименованием; кодировка --- UTF-8 без BOM (критично для парсера HAProxy). Валидация listen-порта и диапазона публикуемых портов выполняется до изменения файлов.

\clearpage
\noindent Пример сгенерированного фрагмента:
\begin{lstlisting}[style=code,language={},caption={Фрагмент listen-секции HAProxy}]
listen chronos_tcp_<идентификатор>
  bind *:<listen-порт>
  mode tcp
  option tcplog
  server backend_<идентификатор> host.docker.internal:<порт> check
\end{lstlisting}

\subsection{Алгоритм реестра TCP-маршрутов}

Подсистема на стороне Master реализует цикл: валидация входных данных (диапазон listen-порта, отсутствие конфликта с уже зарегистрированными маршрутами и зарезервированными портами) $\rightarrow$ обновление JSON-реестра в каталоге динамической конфигурации $\rightarrow$ генерация текстового фрагмента \texttt{chronos-tcp.cfg}. Запись на диск выполняется через временный файл и переименование, чтобы внешний процесс HAProxy не прочитал частично записанный файл.

Выделение listen-порта при автоматическом режиме сводится к перебору кандидатов внутри заданного \texttt{MIN..MAX} с пропуском занятых значений; предпочтительный порт (если передан) проверяется первым. При удалении маршрута реестр и конфигурация приводятся в согласованное состояние за одну операцию; UI отображает результат после повторного чтения с сервера.

\subsection{Фоновая репликация томов}

Служба \texttt{VolumeReplicationHostedService} выполняется на лидирующем экземпляре Master (в MVP это единственный экземпляр). Цикл обхода политик: чтение включённых политик из БД $\rightarrow$ при необходимости запрос свободного места на диске целевого агента $\rightarrow$ листинг объектов в S3 по префиксу и prune по \texttt{MaxCopies} $\rightarrow$ принятие решения о новом снимке с учётом интервала и объёма предыдущей передачи $\rightarrow$ HTTP-вызов агента на endpoint снимка тома. Длительная работа вынесена в фон, чтобы не блокировать основной поток обработки API.

\subsection{Проксирование запросов UI к агентам}

Часть страниц интерфейса запрашивает данные напрямую с узла агента (логи, граф сервисов, диагностика). Master выступает обратным прокси: по идентификатору агента из БД выбирается \texttt{baseUrl}, формируется исходящий HTTP-запрос с тем же путём и методом (с ограничениями безопасности), ответ транслируется клиенту. Таким образом, браузер пользователя не должен иметь прямого сетевого доступа ко всем агентам --- достаточно доверия к Master.

\subsection{Сводная таблица групп UI API}

Таблица~\ref{tab:ui-api-groups} не дублирует полный OpenAPI, а задаёт ориентиры для чтения кода \texttt{UiV1Extensions} и связанных обработчиков.

{\small
\begin{longtable}{@{}p{0.28\linewidth}p{0.66\linewidth}@{}}
\caption{Группы endpoint-ов UI API Master (обзор)} \label{tab:ui-api-groups} \\
\toprule
\textbf{Группа} & \textbf{Назначение} \\
\midrule
\endfirsthead
\multicolumn{2}{c}{\tablename\ \thetable{} --- продолжение} \\
\toprule
\textbf{Группа} & \textbf{Назначение} \\
\midrule
\endhead
Агенты & Список зарегистрированных узлов, детали, действия администратора, связанные с доступностью. \\
Проекты & Обзор размещений, переход к логам и управлению стеком на выбранном агенте. \\
TCP-маршруты & CRUD маршрутов, просмотр сгенерированного фрагмента конфигурации, диагностика конфликтов портов. \\
Архив & Просмотр архивных метаданных проектов, восстановление в работу при поддержке сценария. \\
Песочница & Безопасный предпросмотр Fluent/YAML без записи в продуктивные сущности (при включении функции). \\
Служебные & Здоровье сервиса, версии, вспомогательные данные для отладки UI. \\
\bottomrule
\end{longtable}
}

\subsection{Контейнеризация стенда}

Файл \texttt{deploy/docker/docker-compose.yml} поднимает PostgreSQL (с инициализацией БД), HAProxy со скриптом динамического reload, Loki, Grafana, образы Master и Agent. Такой стенд использовался при функциональных испытаниях.

\subsection{Клиентская часть}

Single Page Application собрана в \texttt{Chronos.Master.Web} (Vite, React, TypeScript, маршрутизация \texttt{react-router-dom}). Оболочка \texttt{ShellLayout} задаёт шапку с логотипом и меню переходов: дашборд, агенты, проекты, архив проектов, TCP routing, sandbox. Запросы к данным агентов выполняются через UI API Master (\texttt{/api/v1/...}), который проксирует вызовы к нужному \texttt{Chronos.Agent}; для части обзорных данных (например, списка агентов) допускаются прямые запросы к маршрутам Master, если они открыты с того же origin.

В подразделе ниже вынесены макеты под скриншоты каждого основного экрана: при подготовке итогового PDF блок с серой рамкой заменяется файлом \texttt{.png}/\texttt{.pdf} (командой \texttt{\textbackslash includegraphics}), размещённым, например, в каталоге \texttt{figures/} рядом с \texttt{main.tex}.

% Скриншоты UI: каталог docs/diplom/фото/ (подпапки по экранам).
% Список проектов: ui.projects.list/1.jpg  |  Карточка проекта: ui.projects/1.jpg, 2.jpg
\subsection{Экраны веб-интерфейса \texttt{Chronos.Master.Web}}

Ниже приведены скриншоты экранов в соответствии с маршрутами React (\texttt{App.tsx}). На рисунки даны отсылки в тексте главы~3 и настоящей главы.

\begin{figure}[!ht]
  \centering
  \IfFileExists{фото/ui.dashboard/1.png}{%
    \includegraphics[width=0.94\linewidth]{ui.dashboard/1.png}%
  }{%
    \IfFileExists{фото/ui.dashboard/1.jpg}{%
      \includegraphics[width=0.94\linewidth]{ui.dashboard/1.jpg}%
    }{%
      \begin{minipage}{0.94\linewidth}\centering
        \fcolorbox{black}{gray!8}{%
          \begin{minipage}{0.92\linewidth}\centering
            \vspace{1.2cm}
            {\small\textit{Скриншот дашборда (\texttt{/ui/dashboard}): положите кадр в \texttt{фото/ui.dashboard/1.png} или \texttt{1.jpg}.}}
            \vspace{1.2cm}
          \end{minipage}}
      \end{minipage}%
    }%
  }
  \caption{Главная панель (дашборд) веб-интерфейса}
  \label{fig:ui-dashboard}
\end{figure}

\begin{figure}[!ht]
  \centering
  \IfFileExists{фото/ui.agents/1.jpg}{%
    \includegraphics[width=0.94\linewidth]{ui.agents/1.jpg}%
  }{%
    \begin{minipage}{0.94\linewidth}\centering
      \fcolorbox{black}{gray!8}{%
        \begin{minipage}{0.92\linewidth}\centering
          \vspace{1.2cm}
          {\small\textit{Скриншот списка агентов (\texttt{/ui/agents}): положите кадр в \texttt{фото/ui.agents/1.jpg}.}}
          \vspace{1.2cm}
        \end{minipage}}
    \end{minipage}%
  }
  \caption{Список агентных узлов}
  \label{fig:ui-agents}
\end{figure}

\begin{figure}[!ht]
  \centering
  \IfFileExists{фото/ui.projects.list/1.jpg}{%
    \includegraphics[width=0.94\linewidth]{ui.projects.list/1.jpg}%
  }{%
    \begin{minipage}{0.94\linewidth}\centering
      \fcolorbox{black}{gray!8}{%
        \begin{minipage}{0.92\linewidth}\centering
          \vspace{1.2cm}
          {\small\textit{Положите скриншот списка проектов в файл \texttt{фото/ui.projects.list/1.jpg} (отдельно от кадров карточки проекта в \texttt{ui.projects/}).}}
          \vspace{1.2cm}
        \end{minipage}}
    \end{minipage}%
  }
  \caption{Список compose-проектов кластера}
  \label{fig:ui-projects}
\end{figure}

\begin{figure}[!ht]
  \centering
  \includegraphics[width=0.94\linewidth]{ui.projects/1.jpg}\\[2mm]
  \includegraphics[width=0.94\linewidth]{ui.projects/2.jpg}
  \caption{Детальная страница проекта}
  \label{fig:ui-project-detail}
\end{figure}

\begin{figure}[!ht]
  \centering
  \includegraphics[width=0.94\linewidth]{ui.projects.archive/1.jpg}
  \caption{Архивированные проекты}
  \label{fig:ui-archived}
\end{figure}

\begin{figure}[!ht]
  \centering
  \includegraphics[width=0.94\linewidth]{ui.routing/1.jpg}\\[2mm]
  \includegraphics[width=0.94\linewidth]{ui.routing/2.jpg}
  \caption{Управление TCP-маршрутизацией (HAProxy)}
  \label{fig:ui-routing}
\end{figure}

\begin{figure}[!ht]
  \centering
  \includegraphics[width=0.94\linewidth]{ui.sandbox/1.jpg}\\[2mm]
  \includegraphics[width=0.94\linewidth]{ui.sandbox/2.jpg}\\[2mm]
  \includegraphics[width=0.94\linewidth]{ui.sandbox/3.jpg}
  \caption{Песочница YAML и Fluent API}
  \label{fig:ui-sandbox}
\end{figure}


\subsection{Практическая часть: запуск из кода}

Практическая демонстрация возможностей Chronos.Core заключается в программном создании экземпляра \texttt{ComposeBuilder}, описании сервисов, портов и томов, вызове валидации и методов \texttt{PublishRemoteAsync}/\texttt{StartRemoteAsync} (или локального \texttt{StartAsync}) с указанием URL агента и ключа API.

\textit{В данной версии текста ВКР листинги конкретного кода автора (фрагменты из репозитория или учебного примера) вставляются в этот подраздел вручную по согласованию с научным руководителем, чтобы сохранить актуальность строк и номеров файлов.} Рекомендуется включить:
\begin{itemize}[leftmargin=1.5cm]
  \item минимальный пример построения одного сервиса с публикацией порта;
  \item пример добавления артефакта деплоя;
  \item пример вызова удалённого старта после \texttt{PublishAsync}.
\end{itemize}

\subsection{Безопасность}

Реализованы проверки API-ключей, ограничения на опасные команды, редактирование потенциально чувствительных подстрок в логах. Для production-развёртывания рекомендуется дополнительно вынести секреты в внешние хранилища и ограничить сетевой доступ к \texttt{docker.sock} на хосте агента.
